% 施設アクセスガイド
\documentclass[12pt,a4paper]{article}

\usepackage{fontspec}
\usepackage{xeCJK}
\setCJKmainfont{Hiragino Mincho ProN}
\setCJKsansfont{Hiragino Kaku Gothic ProN}
\usepackage[margin=2.5cm]{geometry}
\usepackage{amsmath,amssymb}
\usepackage{hyperref}
\usepackage{tcolorbox}
\usepackage{booktabs}
\usepackage{fancyhdr}

\tcbuselibrary{breakable,skins}

\newtcolorbox{facbox}[1][]{colback=blue!5,colframe=blue!60!black,
  fonttitle=\bfseries\sffamily,title={#1},breakable,arc=3mm}
\newtcolorbox{stepbox}[1][]{colback=green!5,colframe=green!50!black,
  fonttitle=\bfseries\sffamily,title={#1},breakable,arc=3mm}
\newtcolorbox{costbox}[1][]{colback=yellow!8,colframe=orange!80!black,
  fonttitle=\bfseries\sffamily,title={#1},breakable,arc=3mm}
\newtcolorbox{honestbox}[1][]{colback=gray!8,colframe=gray!50!black,
  fonttitle=\bfseries\sffamily,title={#1},breakable,arc=3mm}

\setlength{\parskip}{0.8em}
\setlength{\parindent}{0pt}

\pagestyle{fancy}
\fancyhf{}
\rhead{Wright Brothers}
\lhead{施設アクセスガイド}
\cfoot{\thepage}

\begin{document}

\begin{center}
{\Huge\bfseries 施設アクセスガイド}\\[0.5cm]
{\Large $\lambda/4$ vs $\lambda/2$ Lamb シフト実験}\\[0.3cm]
{\Large 世界のオープンアクセス量子施設}\\[1cm]
{\large Wright Brothers Research Program, 2026年4月}
\end{center}

\vspace{0.5cm}
\tableofcontents
\newpage

%====================================================================
\part{実験に必要なもの}
%====================================================================

\section{3つの要素}

\begin{center}
\renewcommand{\arraystretch}{1.5}
\begin{tabular}{clll}
\toprule
\# & 必要なもの & 自作可能？ & 施設で必要 \\
\midrule
1 & 超伝導共振器（$\lambda/4$, $\lambda/2$） & \textbf{✓}（同軸ケーブルを切るだけ） & 不要 \\
2 & トランスモン量子ビット & ✗（ナノ加工が必要） & \textbf{クリーンルーム} \\
3 & 極低温測定 & ✗（希釈冷凍機が必要） & \textbf{低温ラボ} \\
\bottomrule
\end{tabular}
\end{center}

共振器は自作可能（アルミ同軸ケーブル \textyen 5,000）。問題はトランスモンの製作と極低温測定。この2つを提供する施設が必要。

\section{理想的な施設の条件}

\begin{itemize}
\item ナノ加工設備（EBL + 蒸着: トランスモン製作用）
\item 希釈冷凍機（$\leq 20$ mK）
\item マイクロ波測定系（VNA, 低温アンプ, シールドルーム）
\item \textbf{外部ユーザーを受け入れる制度}がある
\item \textbf{個人（機関所属なし）でも利用可能}
\end{itemize}

%====================================================================
\part{世界のオープンアクセス施設}
%====================================================================

%────────────────────────────────────────
\section{★★★ OtaNano（フィンランド）— 最有力}
%────────────────────────────────────────

\begin{facbox}[OtaNano Research Infrastructure]
\textbf{場所}: Aalto University, Espoo, Finland（ヘルシンキ近郊）\\
\textbf{URL}: \url{https://otanano.fi}

\textbf{設備}:
\begin{itemize}
\item \textbf{Micronova クリーンルーム}: EBL, 蒸着装置, スパッタ — トランスモン製作可能
\item \textbf{BlueFors 希釈冷凍機}: 10 mK 以下 — 量子ビット測定可能
\item \textbf{低温マイクロ波測定系}: VNA, HEMT アンプ, 低温ケーブル
\end{itemize}

\textbf{アクセスモデル}: \textbf{Fee-for-service（料金制）}。外部ユーザー歓迎。企業、個人研究者、海外研究者を受け入れ。オンラインでプロポーザルを提出。

\textbf{個人利用}: \textbf{可能}。機関所属は必須ではない（要確認）。料金を払えば使える。

\textbf{言語}: 英語（フィンランドでは全員英語を話す）

\textbf{なぜ最有力か}: ナノ加工 + 希釈冷凍機 + 測定系が\textbf{同一施設内}に揃う世界でも数少ない施設。別の場所にサンプルを運ぶ必要がない。VTT（フィンランド技術研究センター）が隣接しており、商用トランスモン製作サービスもある。
\end{facbox}

%────────────────────────────────────────
\section{★★★ Argonne CNM（アメリカ）— 無料の可能性}
%────────────────────────────────────────

\begin{facbox}[Center for Nanoscale Materials (CNM)]
\textbf{場所}: Argonne National Laboratory, Lemont, IL, USA（シカゴ近郊）\\
\textbf{URL}: \url{https://cnm.anl.gov}

\textbf{設備}:
\begin{itemize}
\item \textbf{クリーンルーム}: EBL（JEOL 高分解能）, 蒸着, エッチング
\item 希釈冷凍機: 隣接する量子プログラムにあるが、ユーザー施設としてのアクセスは要確認
\end{itemize}

\textbf{アクセスモデル}: \textbf{DOE ユーザー施設}。\textbf{非営利・公開研究なら無料}。プロポーザル審査制。

\textbf{個人利用}: DOE ユーザー施設は法的に「いかなる研究者」からのプロポーザルも受け付ける義務がある。機関所属なしでも\textbf{原則として申請可能}（ただし実績がないと審査で不利）。

\textbf{注意}: ナノ加工は強いが、希釈冷凍機の外部ユーザーアクセスは限定的。トランスモンを CNM で作り、測定は別の場所（共同研究先 or OtaNano）で行う二段構えが現実的。

\textbf{コスト}: 公開研究なら\textbf{クリーンルーム使用料ゼロ}。渡航費・滞在費は自己負担。
\end{facbox}

%────────────────────────────────────────
\section{★★ Brookhaven CFN（アメリカ）}
%────────────────────────────────────────

\begin{facbox}[Center for Functional Nanomaterials (CFN)]
\textbf{場所}: Brookhaven National Laboratory, Upton, NY, USA（ニューヨーク郊外）\\
\textbf{URL}: \url{https://cfn.bnl.gov}

Argonne CNM と同じ DOE ユーザー施設モデル。EBL, 蒸着, エッチングが利用可能。公開研究なら無料。希釈冷凍機は隣接プログラムに依存。
\end{facbox}

%────────────────────────────────────────
\section{★★ Cornell CNF（アメリカ）}
%────────────────────────────────────────

\begin{facbox}[Cornell NanoScale Science and Technology Facility]
\textbf{場所}: Cornell University, Ithaca, NY, USA\\
\textbf{URL}: \url{https://cnf.cornell.edu}

\textbf{設備}: 高品質 EBL（JEOL 9500）, 蒸着, エッチング。超伝導デバイス製作の実績多数。

\textbf{アクセス}: プロポーザル制。学術利用は補助金付きレート。外部ユーザー歓迎。

\textbf{注意}: ナノ加工のみ。希釈冷凍機なし。測定は別途。

\textbf{コスト}: クリーンルーム $\sim$\$150--300/日（外部学術レート）。
\end{facbox}

%────────────────────────────────────────
\section{★★ Myfab（スウェーデン）}
%────────────────────────────────────────

\begin{facbox}[Myfab National Nanofabrication Infrastructure]
\textbf{場所}: Chalmers University, Gothenburg, Sweden\\
\textbf{URL}: \url{https://myfab.se}

スウェーデンの国立ナノファブ施設。Chalmers は超伝導量子ビットの世界的拠点（WACQT プロジェクト）。

\textbf{設備}: EBL, 蒸着, エッチング。Chalmers の量子デバイスラボが隣接。

\textbf{アクセス}: 時間制料金。外部ユーザー（海外含む）歓迎。

\textbf{注意}: 希釈冷凍機は WACQT ラボにあるが、外部ユーザーへの開放は要交渉。

\textbf{コスト}: クリーンルーム $\sim$€200--400/日。
\end{facbox}

%====================================================================
\part{推奨ルート}
%====================================================================

\section{Option 1: OtaNano 一括（最もシンプル）}

\begin{stepbox}[全工程を1つの施設で]
\begin{enumerate}
\item OtaNano にプロポーザルを提出（オンライン）
\item 渡航（日本 → ヘルシンキ、フィンランド）
\item Micronova クリーンルームでトランスモン製作（1〜2週間）
\item 同じ施設の希釈冷凍機で $\lambda/4$ vs $\lambda/2$ 測定（1〜2週間）
\item データを持って帰国。論文執筆。
\end{enumerate}
滞在期間: 約1ヶ月。サンプルの輸送不要。
\end{stepbox}

\section{Option 2: Argonne + 測定先（最安）}

\begin{stepbox}[製作無料 + 測定は別途]
\begin{enumerate}
\item Argonne CNM にプロポーザルを提出（公開研究、無料）
\item 渡航（日本 → シカゴ）
\item CNM クリーンルームでトランスモン製作（1〜2週間、\textbf{施設利用料ゼロ}）
\item サンプルを持って OtaNano or 共同研究先に移動
\item 希釈冷凍機で測定
\end{enumerate}
製作費ゼロだが、2箇所に渡航が必要。
\end{stepbox}

\section{Option 3: 商用トランスモン購入 + 測定のみ}

\begin{stepbox}[製作をスキップ]
\begin{enumerate}
\item VTT（フィンランド）or 専門業者からトランスモンチップを購入
\item OtaNano の希釈冷凍機で測定のみ
\end{enumerate}
最も短期間（1〜2週間）。ただしトランスモンの価格が高い。
\end{stepbox}

%====================================================================
\part{必要な訓練}
%====================================================================

\section{事前に身につけるべきスキル}

\begin{stepbox}[クリーンルーム訓練（トランスモン製作用）]
\begin{itemize}
\item \textbf{EBL（電子ビームリソグラフィ）}: パターン設計（KLayout, 無料）、レジスト塗布、描画、現像。各施設でトレーニングコースあり（通常1〜2日）。
\item \textbf{蒸着}: 電子ビーム蒸着装置の操作。角度蒸着（ダリッジブリッジ法）。トレーニング1日。
\item \textbf{酸化}: ジョセフソン接合のトンネルバリア形成。酸化条件の最適化（数回のテスト）。
\item \textbf{リフトオフ}: レジスト除去。
\end{itemize}

多くのユーザー施設は初めてのユーザーに対して\textbf{ハンズオントレーニング}を提供する（施設利用料に含まれる場合が多い）。

\textbf{所要時間}: 初回ユーザーは1週間で基本操作を習得。良いデバイスを安定して作れるまでに2〜4週間。
\end{stepbox}

\begin{stepbox}[低温測定訓練]
\begin{itemize}
\item \textbf{希釈冷凍機の操作}: サンプルマウント、冷却手順、昇温手順。施設スタッフが補助。
\item \textbf{マイクロ波測定}: VNA の使い方、$S_{21}$ 測定、量子ビット分光。
\item \textbf{データ解析}: Python + matplotlib で Lamb シフトを抽出。
\end{itemize}

\textbf{所要時間}: 施設スタッフの補助があれば1週間で最初のデータが取れる。
\end{stepbox}

\begin{honestbox}[事前にオンラインで学べること]
\begin{itemize}
\item KLayout でのマスク設計: \url{https://www.klayout.de/}（無料）
\item circuit QED の基礎: Blais et al. (2004) の論文を読む
\item Python でのデータ解析: numpy, scipy, matplotlib
\item nanoVNA で $\lambda/4$ vs $\lambda/2$ の室温データを取る（Phase 0, \textyen 8,500）
\end{itemize}

\textbf{現地に行く前にこれらを完了しておくと、施設での時間を最大限活用できる。}
\end{honestbox}

%====================================================================
\part{費用の見積もり}
%====================================================================

\section{Option 1: OtaNano 一括}

\begin{costbox}[OtaNano 1ヶ月滞在]
\begin{center}
\renewcommand{\arraystretch}{1.3}
\begin{tabular}{lr}
\toprule
項目 & 費用（概算） \\
\midrule
\multicolumn{2}{l}{\textit{渡航・滞在}} \\
航空券（東京↔ヘルシンキ往復） & \textyen 150,000 \\
宿泊（1ヶ月、Airbnb or ゲストハウス） & \textyen 200,000 \\
食費・交通費（1ヶ月） & \textyen 100,000 \\
\midrule
\multicolumn{2}{l}{\textit{施設利用}} \\
クリーンルーム使用料（2週間） & \textyen 200,000〜400,000 \\
希釈冷凍機使用料（2週間） & \textyen 200,000〜500,000 \\
材料費（ウエハー、蒸着材料、レジスト） & \textyen 50,000 \\
$\lambda/4$, $\lambda/2$ 共振器材料 & \textyen 10,000 \\
\midrule
\textbf{合計} & \textbf{\textyen 910,000〜1,410,000} \\
\bottomrule
\end{tabular}
\end{center}

\textbf{約100万〜140万円}。1ヶ月のフィンランド滞在 + 全実験。
\end{costbox}

\section{Option 2: Argonne（製作無料）+ OtaNano（測定）}

\begin{costbox}[2施設利用]
\begin{center}
\renewcommand{\arraystretch}{1.3}
\begin{tabular}{lr}
\toprule
項目 & 費用（概算） \\
\midrule
\multicolumn{2}{l}{\textit{Argonne（2週間）}} \\
航空券（東京↔シカゴ往復） & \textyen 150,000 \\
宿泊（2週間） & \textyen 100,000 \\
食費・交通費 & \textyen 50,000 \\
クリーンルーム使用料 & \textbf{\textyen 0}（公開研究） \\
材料費 & \textyen 30,000 \\
\midrule
\multicolumn{2}{l}{\textit{OtaNano（2週間）}} \\
航空券（シカゴ↔ヘルシンキ） & \textyen 100,000 \\
宿泊（2週間） & \textyen 100,000 \\
食費・交通費 & \textyen 50,000 \\
希釈冷凍機使用料 & \textyen 200,000〜500,000 \\
\midrule
\textbf{合計} & \textbf{\textyen 780,000〜1,080,000} \\
\bottomrule
\end{tabular}
\end{center}

製作費が無料なので若干安い。ただし2箇所への渡航が必要。
\end{costbox}

\section{Option 3: 商用トランスモン + OtaNano 測定のみ}

\begin{costbox}[最短ルート]
\begin{center}
\renewcommand{\arraystretch}{1.3}
\begin{tabular}{lr}
\toprule
項目 & 費用（概算） \\
\midrule
トランスモンチップ購入（VTT or 業者） & \textyen 100,000〜300,000 \\
航空券（東京↔ヘルシンキ往復） & \textyen 150,000 \\
宿泊（2週間） & \textyen 100,000 \\
食費・交通費 & \textyen 50,000 \\
希釈冷凍機使用料（2週間） & \textyen 200,000〜500,000 \\
共振器材料 & \textyen 10,000 \\
\midrule
\textbf{合計} & \textbf{\textyen 610,000〜1,110,000} \\
\bottomrule
\end{tabular}
\end{center}

クリーンルーム不要。最短2週間。ただしトランスモンの調達に時間がかかる可能性。
\end{costbox}

%====================================================================
\part{タイムライン}
%====================================================================

\begin{center}
\renewcommand{\arraystretch}{1.5}
\begin{tabular}{clp{6cm}}
\toprule
期間 & フェーズ & 内容 \\
\midrule
\textit{渡航前} & 準備（1〜3ヶ月） & Phase 0 完了（nanoVNA）、KLayout 習得、cQED 論文を arXiv 投稿、OtaNano にプロポーザル提出 \\
\textit{Week 1-2} & 製作 & クリーンルームでトランスモン製作。同時に共振器を準備 \\
\textit{Week 3} & セットアップ & 冷凍機にサンプルマウント。配線。冷却開始 \\
\textit{Week 4} & 測定 & $\lambda/2$ と $\lambda/4$ の Lamb シフトを測定。データ取得 \\
\textit{帰国後} & 解析・執筆 & データ解析。cQED 論文（結果報告）+ Boyer 論文（数論的解釈）を投稿 \\
\bottomrule
\end{tabular}
\end{center}

%====================================================================
\part{購入リスト}
%====================================================================

\section{渡航前に日本で購入}

\begin{center}
\renewcommand{\arraystretch}{1.3}
\begin{tabular}{lrl}
\toprule
品目 & 価格 & 備考 \\
\midrule
nanoVNA & \textyen 3,000 & Phase 0 用 \\
セミリジッド同軸ケーブル（Al or Cu）& \textyen 5,000 & 共振器材料 \\
SMA コネクタ各種 & \textyen 2,000 & \\
SMA ショートキャップ & \textyen 500 & $\lambda/2$ 用 \\
ハンダごて + ハンダ & \textyen 3,000 & 持っていれば不要 \\
\midrule
\textbf{合計} & \textbf{\textyen 13,500} & \\
\bottomrule
\end{tabular}
\end{center}

\section{現地で調達 or 施設で使用}

\begin{center}
\renewcommand{\arraystretch}{1.3}
\begin{tabular}{ll}
\toprule
品目 & 入手方法 \\
\midrule
Si ウエハー（トランスモン基板） & 施設の在庫 or 現地購入 \\
PMMA レジスト & 施設の在庫 \\
Al 蒸着用ワイヤー & 施設の在庫 \\
低温同軸ケーブル & 施設の在庫 or レンタル \\
SMA アダプタ・アッテネータ & 施設の在庫 \\
\bottomrule
\end{tabular}
\end{center}

%====================================================================
\part{正直な評価}
%====================================================================

\begin{honestbox}[リスク]
\textbf{最大のリスク: トランスモンの製作に失敗する}

ジョセフソン接合の酸化条件は微妙で、初めての製作では9割失敗する。5〜10回のイテレーションが必要。2週間で成功する保証はない。

\textbf{対策}:
\begin{itemize}
\item Option 3（商用トランスモン購入）で製作リスクを回避
\item 施設のスタッフに「トランスモンを作りたい」と事前に伝え、レシピを共有してもらう
\item 事前に論文を読んで酸化条件のパラメータ範囲を把握
\end{itemize}
\end{honestbox}

\begin{honestbox}[プロポーザルのコツ]
施設は「何をしたいか」「なぜその施設が必要か」「成果をどう公開するか」を聞く。

書くべきこと:
\begin{itemize}
\item 「$\lambda/4$ と $\lambda/2$ 共振器における多モード真空揺らぎの量子ビット周波数安定性への影響を測定する」
\item 「結果は査読付き論文として公開する」（これで DOE 施設は無料になる）
\item 「所要期間: 4週間、クリーンルーム2週間 + 低温測定2週間」
\end{itemize}

\textbf{書かないこと}: ζ関数、素数、ワープドライブ。
\end{honestbox}

\begin{honestbox}[施設利用料は「概算」]
上記の金額は公開情報と一般的なレートからの推定。実際の料金は施設に直接問い合わせる必要がある。OtaNano は \url{https://otanano.fi} から見積もり依頼が可能。
\end{honestbox}

\section{推奨}

\begin{costbox}[結論]
\textbf{最もシンプル}: Option 1（OtaNano 一括）。1箇所で全て完結。約100万〜140万円。

\textbf{最も安い}: Option 2（Argonne + OtaNano）。製作無料。約80万〜110万円。ただし2箇所渡航。

\textbf{最も速い}: Option 3（商用トランスモン + OtaNano 測定）。2週間。約60万〜110万円。

\textbf{共通}: フィンランドの OtaNano が測定の拠点。ナノ加工 + 希釈冷凍機 + 測定系が揃う世界有数のオープンアクセス施設。
\end{costbox}

\vfill

\begin{center}
\rule{10cm}{0.4pt}\\[0.5cm]
{\Large\bfseries 装置は世界中にある。}\\[0.2cm]
{\Large\bfseries 足りないのはパスポートと航空券だけ。}\\[0.3cm]
{\small Wright Brothers Research Program, 2026}
\end{center}

\end{document}
